\documentclass[a4paper,12pt]{article}

\usepackage[T2A]{fontenc}
\usepackage[utf8]{inputenc}
\usepackage[russian]{babel}

\usepackage{amsmath,amssymb,amsthm,mathtools}
\usepackage{helvet}
\renewcommand{\familydefault}{\sfdefault}
\usepackage{icomma}

\usepackage[a4paper,margin=2.2cm]{geometry}
\usepackage{microtype}

\usepackage{tikz}
\usetikzlibrary{calc}
\usetikzlibrary{arrows.meta,positioning,shapes.geometric}

\usepackage{tkz-euclide}

\usepackage{pgfplots}
\pgfplotsset{compat=1.18}

\usepackage{tabularx,booktabs,longtable}
\usepackage{enumitem}
\usepackage{hyperref}
\usepackage{parskip}
\usepackage{fancyhdr}
\usepackage{setspace}
\usepackage{xcolor}

\definecolor{accent}{HTML}{008080}
\definecolor{darkaccent}{HTML}{004D4D}
\definecolor{textgray}{HTML}{333333}

\hypersetup{
  colorlinks=true,
  linkcolor=darkaccent,
  urlcolor=darkaccent
}

\geometry{headheight=25pt}

\onehalfspacing
\setlength{\parindent}{0pt}

\setlist[itemize]{
  leftmargin=1.6em,
  itemsep=0.25em,
  topsep=0.3em
}

\setlist[enumerate]{
  leftmargin=1.8em,
  itemsep=0.45em,
  topsep=0.3em
}

\renewcommand{\arraystretch}{1.3}

\pagestyle{fancy}
\fancyhf{}
\fancyhead[L]{\small\textcolor{textgray}{ЕГЭ. Степени и корни}}
\fancyhead[R]{\small\textcolor{textgray}{Екатерина Гнедкова}}
\fancyfoot[C]{\thepage}

\renewcommand{\headrulewidth}{0.3pt}
\renewcommand{\headrule}{
  \hbox to\headwidth{
    \color{accent}
    \leaders\hrule height \headrulewidth\hfill
  }
}

\newcommand{\sectionline}{
  \par
  \vspace{0.2em}
  {\color{accent}\hrule height 0.6pt}
  \vspace{0.8em}
}

\newcommand{\checkitem}[1]{\item[$\square$] #1}

\tikzset{
  flowstart/.style={
    ellipse,
    draw=darkaccent,
    line width=0.9pt,
    minimum width=3.4cm,
    minimum height=1.0cm,
    align=center,
    font=\small
  },
  flowact/.style={
    rounded corners=3pt,
    draw=darkaccent,
    line width=0.9pt,
    text width=6.0cm,
    minimum height=1.05cm,
    align=center,
    font=\small,
    inner sep=6pt
  },
  flowcond/.style={
    diamond,
    aspect=2.0,
    draw=darkaccent,
    line width=0.9pt,
    text width=3.6cm,
    align=center,
    font=\small,
    inner sep=2pt
  },
  flowarrow/.style={
    -{Latex[length=3mm,width=2mm]},
    line width=0.9pt,
    draw=darkaccent
  },
  flowlabel/.style={
    font=\small,
    fill=white,
    inner sep=1.5pt
  }
}

\begin{document}

% =========================================================
% ТИТУЛЬНЫЙ ЛИСТ
% =========================================================

\begin{titlepage}
\thispagestyle{empty}
\centering

\vspace*{2.1cm}

{\Large\textcolor{darkaccent}{\textbf{Чек-лист}}\par}

\vspace{0.7cm}

{\huge\bfseries
Степени, корни и приведение\\[0.15cm]
к единому степенному виду
\par}

\vspace{0.9cm}

{\large
Подготовка к ЕГЭ: вычисления и преобразования
\par}

\vspace{1.8cm}

{\large Екатерина Гнедкова\par}

\vspace{0.45cm}

{\normalsize
Лёвин Артём Александрович эксклюзивно для Екатерины Гнедковой
\par}

\vspace{0.45cm}

{\normalsize \today\par}

\vfill

\begin{tikzpicture}[remember picture,overlay]
\draw[
  accent,
  line width=1.0pt
]
  ([xshift=1.8cm,yshift=1.9cm]current page.south west)
  .. controls +(3.1cm,1.2cm) and +(-3.1cm,-0.6cm) ..
  ([xshift=-1.8cm,yshift=2.3cm]current page.south east);
\end{tikzpicture}

\end{titlepage}

% =========================================================
% ОСНОВНОЙ ЧЕК-ЛИСТ
% =========================================================

\section*{1. Главная идея темы}
\sectionline

В вычислительных задачах со степенями цель обычно состоит в том, чтобы
\textbf{привести основания к удобному единому виду}, а затем применить
свойства степеней.

Наиболее часто преобразуются:

\begin{itemize}
  \item составные числа:
  \[
  25=5^2,\qquad
  32=2^5,\qquad
  81=3^4,\qquad
  144=2^4\cdot3^2;
  \]

  \item корни:
  \[
  \sqrt[q]{a^p}=a^{p/q}
  \]
  при допустимых значениях;

  \item дроби и отрицательные степени:
  \[
  \frac{1}{a^n}=a^{-n},\qquad a\ne0;
  \]

  \item десятичные показатели:
  \[
  0{,}6=\frac35,
  \qquad
  1{,}25=\frac54.
  \]
\end{itemize}

\subsection*{Что важно различать}

Преобразование в степенной вид обычно относится к
\textbf{основанию}.

Если в показателе степени присутствует, например, $\sqrt7$, его без
необходимости не заменяют на $7^{1/2}$: такое преобразование часто
только усложняет выражение.

% ---------------------------------------------------------

\section*{2. Базовые свойства степеней}
\sectionline

Пусть все записанные выражения определены.

\begin{center}
\begin{tabularx}
{\linewidth}
{@{}>{\bfseries\raggedright\arraybackslash}p{0.29\linewidth}X@{}}

\toprule
Свойство & Формула и условие применения \\
\midrule

Произведение одинаковых оснований
&
\[
a^m\cdot a^n=a^{m+n}.
\]
Между степенями должно быть умножение, основания должны совпадать.
\\

Частное одинаковых оснований
&
\[
\frac{a^m}{a^n}=a^{m-n},
\qquad a\ne0.
\]
\\

Степень произведения
&
\[
(ab)^n=a^n b^n.
\]
В обратную сторону:
\[
a^n b^n=(ab)^n.
\]
\\

Степень частного
&
\[
\left(\frac ab\right)^n=\frac{a^n}{b^n},
\qquad b\ne0.
\]
\\

Степень степени
&
\[
(a^m)^n=a^{mn}.
\]
Поэтому
\[
a^{mn}=(a^m)^n=(a^n)^m.
\]
\\

Нулевая степень
&
\[
a^0=1,
\qquad a\ne0.
\]
\\

Отрицательная степень
&
\[
a^{-n}=\frac{1}{a^n},
\qquad a\ne0.
\]
\\

Дробная степень
&
Для $a>0$:
\[
a^{p/q}=\sqrt[q]{a^p},
\qquad
q\in\mathbb N,\quad q\ge2.
\]
\\

\bottomrule
\end{tabularx}
\end{center}

\subsection*{Правило «читать формулу в обе стороны»}

Равенство полезно воспринимать как инструмент преобразования в обе
стороны. Например,

\[
a^{m+n}=a^m a^n,
\qquad
a^{mn}=(a^m)^n.
\]

Это позволяет специально \textbf{разбивать степень}, если такой ход
создаёт нужное основание или показатель.

% ---------------------------------------------------------

\section*{3. Корни и дробные показатели}
\sectionline

Для положительного основания переход между корнем и степенью выполняется
по правилу

\[
\sqrt[q]{a^p}=a^{p/q}.
\]

Полезное правило для запоминания:

\[
\boxed{
\text{внутренняя степень}
\; \div \;
\text{показатель корня}
}
\]

Примеры:

\[
\sqrt3=3^{1/2},
\qquad
\sqrt[3]{5}=5^{1/3},
\qquad
\sqrt[3]{3^4}=3^{4/3}.
\]

Если индекс квадратного корня не записан, подразумевается число $2$:

\[
\sqrt a=\sqrt[2]{a}=a^{1/2},
\qquad a\ge0.
\]

\subsection*{ОДЗ и ограничения}

\begin{itemize}
  \item Для корня чётной степени
  \[
  \sqrt[2k]{A}
  \]
  необходимо
  \[
  A\ge0.
  \]

  \item Если корень находится в знаменателе, требуется дополнительно
  исключить ноль. Например,
  \[
  \frac1{\sqrt A}
  \]
  определено при
  \[
  A>0.
  \]

  \item В выражениях вида
  \[
  a^{-r}
  \]
  основание должно удовлетворять условию
  \[
  a\ne0.
  \]

  \item Формулу
  \[
  a^{p/q}=\sqrt[q]{a^p}
  \]
  в задачах с произвольными вещественными показателями безопасно
  применять при
  \[
  a>0.
  \]
  Для отрицательных оснований допустимость зависит от конкретного
  рационального показателя.
\end{itemize}

% ---------------------------------------------------------

\section*{4. Составные числа как степени простых множителей}
\sectionline

Составное число имеет более двух натуральных делителей.

В задачах на степени составные числа удобно раскладывать на простые
множители:

\[
144=2^4\cdot3^2,
\qquad
81=3^4,
\qquad
35=5\cdot7,
\qquad
12=2^2\cdot3.
\]

Такое разложение создаёт одинаковые основания и открывает возможность
применять свойства степеней.

\subsection*{Пример}

\[
\frac{35^{-4{,}7}\cdot7^{5{,}7}}{5^{-3{,}7}}
=
\frac{(5\cdot7)^{-4{,}7}\cdot7^{5{,}7}}
{5^{-3{,}7}}.
\]

Используем степень произведения:

\[
=
\frac{5^{-4{,}7}\cdot7^{-4{,}7}\cdot7^{5{,}7}}
{5^{-3{,}7}}.
\]

Группируем одинаковые основания:

\[
=
5^{-4{,}7-(-3{,}7)}
\cdot
7^{-4{,}7+5{,}7}.
\]

Следовательно,

\[
=
5^{-1}\cdot7
=
\frac75.
\]

% ---------------------------------------------------------

\section*{5. Типовой метод решения вычислительных задач}
\sectionline

\begin{enumerate}
  \item Найдите \textbf{преобразуемые основания}:
  составные числа, корни, дроби.

  \item Разложите составные числа и переведите корни и дроби в
  степенной вид.

  \item Стремитесь получить одинаковые основания или одинаковые
  показатели.

  \item Применяйте формулы степеней последовательно, по одному
  преобразованию за шаг.

  \item Аккуратно вычисляйте показатели:
  \begin{itemize}
    \item при умножении одинаковых оснований показатели складываются;
    \item при делении показатели вычитаются;
    \item при возведении степени в степень показатели перемножаются.
  \end{itemize}

  \item Если появился отрицательный показатель, при необходимости
  вернитесь к дроби.

  \item Проверьте ОДЗ и ограничения.

  \item Запишите окончательный ответ.
\end{enumerate}

\subsection*{Пример}

\[
\left(
\frac{2^{1/3}\cdot2^{1/4}}
{\sqrt[12]{2}}
\right)^2.
\]

Переводим корень в степенной вид:

\[
\sqrt[12]{2}=2^{1/12}.
\]

Тогда

\[
\left(
\frac{2^{1/3}\cdot2^{1/4}}
{2^{1/12}}
\right)^2
=
\left(
2^{1/3+1/4-1/12}
\right)^2.
\]

Приводим показатели к общему знаменателю:

\[
\frac13+\frac14-\frac1{12}
=
\frac4{12}+\frac3{12}-\frac1{12}
=
\frac6{12}
=
\frac12.
\]

Следовательно,

\[
\left(2^{1/2}\right)^2=2.
\]

% =========================================================
% ОТДЕЛЬНАЯ СТРАНИЦА БЛОК-СХЕМЫ
% =========================================================

\clearpage
\pagestyle{empty}

\begin{center}

\begin{tikzpicture}[node distance=0.45cm]

\node[
  text=darkaccent,
  font=\Large\bfseries,
  align=center,
  text width=14cm
] (title)
{
Алгоритм: приведение выражения\\
к единому степенному виду
};

\node[
  flowstart,
  below=0.35cm of title
] (start)
{
Получено выражение со степенями
};

\node[
  flowcond,
  below=of start
] (convertq)
{
Есть составные числа,
корни или дроби
в основаниях?
};

\node[
  flowact,
  below=0.58cm of convertq
] (convert)
{
Разложить составные числа;
корни и дроби перевести
в степенной вид
};

\node[
  flowcond,
  below=0.58cm of convert
] (formulaq)
{
Можно применить
свойство степеней
к одинаковым основаниям
или показателям?
};

\node[
  flowact,
  below=0.58cm of formulaq
] (formula)
{
Применить свойство степеней
и аккуратно упростить показатели
};

\node[
  flowcond,
  below=0.58cm of formula
] (doneq)
{
Получено выражение,
которое можно вычислить
или записать компактно?
};

\node[
  flowact,
  right=0.55cm of formulaq,
  text width=3.6cm
] (rearrange)
{
Перегруппировать
множители,
разбить степень
или вынести
общий показатель
};

\node[
  flowact,
  below=0.58cm of doneq
] (odz)
{
Проверить ОДЗ и ограничения
};

\node[
  flowstart,
  below=of odz
] (finish)
{
Записать ответ
};

\draw[flowarrow]
  (start) -- (convertq);

\draw[flowarrow]
  (convertq)
  --
  node[flowlabel,right]{да}
  (convert);

\draw[flowarrow]
  (convert) -- (formulaq);

\draw[flowarrow]
  (formulaq)
  --
  node[flowlabel,right]{да}
  (formula);

\draw[flowarrow]
  (formula) -- (doneq);

\draw[flowarrow]
  (doneq)
  --
  node[flowlabel,right]{да}
  (odz);

\draw[flowarrow]
  (odz) -- (finish);

\draw[flowarrow]
  (convertq.west)
  -- ++(-0.85,0)
  |-
  node[pos=0.24,flowlabel,left]{нет}
  (formulaq.west);

\draw[flowarrow]
  (formulaq.east)
  --
  node[flowlabel,above]{нет}
  (rearrange.west);

\draw[flowarrow]
  (rearrange.south)
  |-
  (formula.east);

\draw[flowarrow]
  (doneq.west)
  -- ++(-1.75,0)
  |-
  node[pos=0.18,flowlabel,left]{нет}
  (convertq.west);

\end{tikzpicture}

\end{center}

\clearpage
\pagestyle{fancy}

% =========================================================

\section*{6. Границы применимости и типичные ошибки}
\sectionline

\begin{itemize}

  \item Нельзя складывать показатели при сложении:

  \[
  a^m+a^n\ne a^{m+n}.
  \]

  \item Нельзя распределять степень по сумме:

  \[
  (a+b)^n\ne a^n+b^n
  \]

  в общем случае.

  \item При делении показатели вычитаются \textbf{целиком}.

  Например,

  \[
  \frac{5^{\alpha}}
  {5^{2\sqrt7-1}}
  =
  5^{\alpha-(2\sqrt7-1)}.
  \]

  Скобки вокруг показателя знаменателя позволяют сохранить знак
  каждого слагаемого.

  \item При подстановке выражения вместо числа используйте скобки:

  \[
  25^{0{,}4}
  =
  (5^2)^{0{,}4}.
  \]

  \item Не путайте разложение на множители со сложением:

  \[
  144=2^4\cdot3^2,
  \]

  но

  \[
  144\ne2^4+3^2.
  \]

  \item Формулы произведения степеней требуют одинаковых оснований:

  \[
  2^m\cdot3^n
  \]

  нельзя заменить выражением вида

  \[
  2^{m+n}.
  \]

  \item Формулы общего показателя требуют одинаковых показателей:

  \[
  a^n b^n=(ab)^n.
  \]

  Если показатели различны, сначала необходимо выполнить другое
  преобразование.

  \item Не преобразуйте показатели степени без необходимости.
  Главная задача в типовых вычислительных примерах ---
  упростить \textbf{основания} и создать возможность применить формулы.

\end{itemize}

% ---------------------------------------------------------

\section*{7. Краткий чек-лист самопроверки}
\sectionline

\begin{itemize}

  \checkitem{
  Я различаю основание и показатель степени.
  }

  \checkitem{
  Я знаю свойства произведения, частного и степени степени.
  }

  \checkitem{
  Я умею читать свойства степеней в обе стороны.
  }

  \checkitem{
  Я умею переводить корень в дробную степень.
  }

  \checkitem{
  Я умею переводить отрицательную степень в дробь и обратно.
  }

  \checkitem{
  Я раскладываю составные числа на простые множители.
  }

  \checkitem{
  Я сначала ищу способ привести основания к единому виду.
  }

  \checkitem{
  Я не применяю свойства произведения к сумме или разности.
  }

  \checkitem{
  Я учитываю ОДЗ корней и знаменателей.
  }

  \checkitem{
  Я проверяю арифметику показателей перед записью ответа.
  }

\end{itemize}

% =========================================================
% ТРЕНИРОВОЧНЫЙ БЛОК
% =========================================================

\clearpage

\section*{Тренировочный блок}
\sectionline

Задания расположены по нарастающей сложности.
Решения оформите самостоятельно, используя общий алгоритм приведения к
степенному виду.

\begin{enumerate}

  % -------------------------------------------------------
  % 3 задания среднего уровня
  % -------------------------------------------------------

  \item Вычислите:

  \[
  2^{3/5}\cdot8^{4/5}.
  \]

  \item Вычислите:

  \[
  \frac{3^{7/4}}
  {\sqrt[4]{3^3}}.
  \]

  \item Вычислите:

  \[
  \left(
  \frac{2^{1/3}\cdot2^{1/6}}
  {\sqrt[6]{2}}
  \right)^3.
  \]

  % -------------------------------------------------------
  % 6 заданий выше среднего
  % -------------------------------------------------------

  \item Упростите:

  \[
  \frac{12^{3/2}}
  {2^{3/2}\cdot3^{3/2}}.
  \]

  \item Вычислите:

  \[
  25^{-0{,}7}\cdot5^{2{,}4}.
  \]

  \item Вычислите:

  \[
  \frac{18^{5/3}}
  {2^{5/3}\cdot3^{10/3}}.
  \]

  \item Вычислите:

  \[
  \left(
  \frac{
    2^{2/3}\cdot\sqrt[6]{4}
  }
  {
    \sqrt[3]{2}
  }
  \right)^3.
  \]

  \item Упростите:

  \[
  \frac{
    7^{\sqrt3+2}\cdot7^{1-\sqrt3}
  }
  {49}.
  \]

  \item Вычислите:

  \[
  \frac{
    35^{-2{,}7}\cdot7^{3{,}7}
  }
  {
    5^{-1{,}7}
  }.
  \]

  % -------------------------------------------------------
  % 1 сложное задание
  % -------------------------------------------------------

  \item Вычислите:

  \[
  \frac{1}{32}
  \left(
  \frac{
    2^{1/3}\cdot8^{5/12}
  }
  {
    \sqrt[4]{2^3}
  }
  \right)^6.
  \]

\end{enumerate}

% =========================================================
% ОТВЕТЫ
% =========================================================

\clearpage

\section*{Ответы к тренировочному блоку}
\sectionline

\begin{table}[ht]
\centering
\caption{Ответы}

\begin{tabularx}{\linewidth}{@{}cX@{}}
\toprule
№ & Ответ \\
\midrule

1 & $8$ \\

2 & $3$ \\

3 & $2$ \\

4 & $2\sqrt2$ \\

5 & $5$ \\

6 & $1$ \\

7 & $4$ \\

8 & $7$ \\

9 & $\dfrac75=1{,}4$ \\

10 & $1$ \\

\bottomrule
\end{tabularx}

\end{table}

\end{document}